
{\actuality} 

Инфракрасное излучение было открыто английским астрономом У.~Гершелем в 1800 году, и вскоре после этого были обнаружены такие его свойства, как способность поглощаться с выделением тепла и существование окон прозрачности атмосферы в некоторых спектральных диапазонах. Однако генерация когерентного излучения среднего ИК-диапазона (3--50 мкм) стала возможна только после изобретения лазеров и развития нелинейной оптики. В настоящее время для этих целей широко используются мощные одномодовые волоконные лазеры, которые благодаря сочетанию высокой пиковой (до 100~кВт) и средней (до 10~кВт) оптической мощности, высокому качеству пучка ($M^2 < 1.3$) и коэффициенту полезного действия (более 45\% от розетки) успешно конкурируют с твердотельными аналогами. Для преобразования их излучения в средний ИК--диапазон применяются нелинейно-оптические (н-о) кристаллы, позволяющие с помощью эффектов второго порядка (генерация разностной, суммарной частот) расширять доступный спектральный диапазон от ультрафиолета до среднего ИК. Особое место занимает область длин волн около $3$~мкм, поскольку в неё попадают пики поглощения воды, молекулярных связей, а также несколько окон прозрачности атмосферы. Поэтому излучение в данной спектральной области широко используется в медицине (дерматология, хирургия), спектроскопии, дистанционном зондировании атмосферы, лазерной связи в свободном пространстве~\cite{jean2003mid, haas2016advances, lambert2015differential, grillot2025progress}. Для генерации излучения в диапазоне 3 мкм используются различные твердотельные активные среды, включая кристалл Er:YAG, фторидное волокно Ho:ZBLAN, а также халькогениды (ZnS, ZnSe), легированные ионами Cr²⁺ или Fe²⁺.
%Такие лазеры активно применяется в коммерческих медицинских системах~\cite{bader2006indications, jean2003mid}.
Альтернативным методом, позволяющим создавать компактные и эргономичные устройства с волоконным инструментом, обеспечивающим возможность перестройки длины волны и узкую линию излучения, является генерация разностной частоты (ГРЧ) в кристалле ниобата лития с регулярной доменной структурой (PPLN) при накачке двумя импульсными волоконными лазерами, излучающими на длинах волн 1030~нм и 1560~нм, импульсы которых синхронизированы по времени и распространяются по одному оптическому волокну~\cite{larionov2021pp}.

Однако при создании стабильных и эффективных источников среднего ИК-диапазона следует учитывать две группы явлений, связанных с н-о и тепловыми процессами в оптических волокнах (ОВ) и н-о кристаллах.

Первая группа связана с синхронным распространением мощных оптических импульсов на двух длинах волн по ОВ, которое используется для доставки излучения к н-о кристаллу. С повышением пиковой мощности в ОВ проявляются н-о эффекты. Явления, обусловленные нелинейной поляризуемостью третьего порядка, такие как четырёхволновое смешение (ЧВС), приводят к искажению спектра и снижению мощности в исходных спектральных компонентах~\cite{agrawal2019}. В ОВ, не являющихся одномодовыми хотя бы на одной из длин волн излучения, картина ЧВС усложняется. В таких условиях возможно одновременное протекание как одномодового, так и межмодового ЧВС, поэтому необходимо выяснить, какие именно типы ЧВС реализуются в конкретной ситуации и почему они влияют друг на друга.


Помимо эффектов третьего порядка, в ОВ, несмотря на его центросимметричность, возможна реализация оптически индуцированных эффектов второго порядка, в частности, генерации второй гармоники (ГВГ)~\cite{osterberg1986dye}. Этот эффект снижает мощность излучения основной частоты, ухудшает качество пучка и уменьшает порог оптического разрушения нелинейно-оптических кристаллов. Механизм фотоиндуцированной ГВГ связан с формированием объёмной решётки квадратичной нелинейности при длительном облучении ОВ. Эффективность этого процесса определяется условиями фазового синхронизма, поэтому исследование эволюции соответствующих кривых фазового синхронизма (КФС) в процессе записи решётки может дать важную информацию о динамике её формирования и, в конечном итоге, о её микроскопической природе. Однако до настоящего времени такие экспериментальные исследования, особенно в отсутствие затравочного излучения второй гармоники (что типично для практических лазерных систем, где ГВГ является нежелательным эффектом), проведены не были.

Вторая группа явлений возникает при н-о преобразовании в н-о кристалле. Любой реальный кристалл обладает хоть и малым, но ненулевым коэффициентом оптического поглощения, причём на различных длинах волн излучения коэффициент также будет различаться. Например, для PPLN на длине волны 3 мкм коэффициент поглощения составляет порядка $10^{-2}~\text{см}^{-1}$, в то время как на длинах волн $1.03$~мкм и $1.56$~мкм он составляет $10^{-3}~\text{см}^{-1}$ и $10^{-4}~\text{см}^{-1}$ соответственно.~\cite{leidinger2015comparative} Оптическое поглощение лазерного излучения приводит к разогреву кристалла, изменению показателя преломления и нарушению условий фазового синхронизма, что снижает эффективность н-о преобразования. Для учёта и компенсации тепловых эффектов, таких, как индуцированная фазовая расстройка или термолинзирование, необходимо точное знание коэффициента поглощения. Перспективным методом его измерения является пьезоэлектрическая резонансная лазерная калориметрия (ПРЛК), позволяющая определять усреднённую температуру пьезоэлектрического кристалла по сдвигу частоты одной его собственных акустических мод, бесконтактно возбуждаемой внешним переменным электрическим полем за счёт обратного пьезоэлектрического эффекта~\cite{aloyan2022acoustic}. Для обработки температурных кинетик применяют три стандартных подхода: экспоненциальную методику, в которой коэффициент оптического поглощения определяют по кинетике нагрева; импульсную методику, основанную на анализе кинетики охлаждения; а также градиентную методику, основанную на сопоставлении скоростей изменения температуры при нагреве и остывании для одинакового разогрева. Общим недостатком этих подходов является то, что они не учитывают нестационарность температуры окружающей среды и коэффициента теплоотдачи, что приводит к существенному разбросу результатов.

Настоящая диссертационная работа посвящена исследованию н-о процессов при распространении мощных оптических импульсов по ОВ, а также тепловых процессов в н-о кристаллах, ограничивающих эффективность и стабильность гибридного источника излучения на длине волны 3~мкм.

\ifsynopsis
%Этот абзац появляется только в~автореферате.
%Для формирования блоков, которые будут обрабатываться только в~автореферате,
%заведена проверка условия \verb!\!\verb!ifsynopsis!.
%Значение условия задаётся в~основном файле документа (\verb!synopsis.tex! для
%автореферата).
\else
%Этот абзац появляется только в~диссертации.
%Через проверку условия \verb!\!\verb!ifsynopsis!, задаваемого в~основном файле
%документа (\verb!dissertation.tex! для диссертации), можно сделать новую
%команду, обеспечивающую появление цитаты в~диссертации, но~не~в~автореферате.
\fi

% {\progress}
% Этот раздел должен быть отдельным структурным элементом по
% ГОСТ, но он, как правило, включается в описание актуальности
% темы. Нужен он отдельным структурынм элемементом или нет ---
% смотрите другие диссертации вашего совета, скорее всего не нужен.

Таким образом, {\aim} диссертационной работы является исследование н-о эффектов второго и третьего порядка, развивающихся в ОВ доставки излучения, а также разработка подходов к повышению точности измерения коэффициента оптического поглощения в н-о кристаллах для источника излучения на длине волны вблизи 3~мкм.

Для~достижения поставленной цели необходимо было решить следующие {\tasks}:
\begin{enumerate}[beginpenalty=10000] % https://tex.stackexchange.com/a/476052/104425
  \item выяснить причины взаимного влияния процессов четырёхволнового смешения при распространении синхронизированных оптических импульсов двухволнового лазерного источника в маломодовом ОВ и разработать математическую модель, описывающую это взаимодействие;
  
  \item исследовать эволюцию кривых фазового синхронизма при фотоиндуцированной генерации второй гармоники в оптическом волокне в процессе формирования решётки нелинейной поляризуемости второго порядка и сопоставить экспериментальные результаты с модельными;

  \item проанализировать основные источники погрешности измерения коэффициента оптического поглощения в нелинейно-оптических кристаллах методом ПРЛК и предложить способы их компенсации.
\end{enumerate}


{\novelty}
\begin{enumerate}[beginpenalty=10000] % https://tex.stackexchange.com/a/476052/104425
  \item Установлено, что взаимное влияние одномодового и межмодового четырёхволнового смешения в маломодовом оптическом волокне обусловлено наличием у этих н-о процессов общей антистоксовой компоненты. Предложен метод подавления нежелательных спектральных компонент с помощью модового фильтра.
  
  \item Обнаружено, что в процессе формирования решётки квадратичной нелинейной восприимчивости при фотоиндуцированной генерации второй гармоники в германосиликатном ОВ ширина кривых фазового синхронизма монотонно уменьшается, их пик сдвигается к расстройке, определяемой фазовой самомодуляцией и кросс-модуляцией, при этом форма кривых синхронизма сохраняет выраженную асимметрию.  
  
  \item Для экспоненциального, импульсного и градиентного методов обработки температурных кинетик при разогреве н-о кристалла лазерным излучением предложены поправочные выражения, позволяющие компенсировать нестационарность температуры окружающей среды и коэффициента теплообмена.
  
\end{enumerate}

{\influence}:

\begin{enumerate}[beginpenalty=10000] 
	\item Предложенный метод подавления межмодового четырёхволнового смешения основан на использовании модового фильтра в виде участка одномодового ОВ, согласованного по модовому диаметру с ОВ доставки. Применение метода позволяет увеличить допустимую длину ОВ доставки с 2.5 м до 3.0 м без снижения эффективности генерации разностной частоты, что важно для создания практичных и функциональных лазерных источников среднего ИК-диапазона.
	
	\item Полученные зависимости эволюции кривых фазового синхронизма при оптическом полинге ОВ позволяют определить изменение фазовой расстройки в процессе формирования решётки квадратичной нелинейности. Это создаёт основу для будущих исследований по управлению фазой накачки с целью либо подавления паразитной ГВГ, либо её эффективной генерации.
	
	\item Разработанные поправочные выражения для обработки температурных кинетик н-о кристаллов при разогреве лазерным излучением позволяют снизить разброс измерений коэффициента оптического поглощения в три раза, повышая точность диагностики нелинейно-оптических кристаллов и достоверность оценки их теплового режима при работе в мощных лазерных системах.
	
\end{enumerate}
%
%{\methods} \ldots

{\defpositions}
\begin{enumerate}[beginpenalty=10000] % https://tex.stackexchange.com/a/476052/104425
  \item В маломодовом оптическом волокне взаимное влияние одномодового и межмодового четырёхволнового смешения синхронизированных оптических импульсов двухволнового лазерного источника обусловлено перекрытием спектров усиления их общих антистоксовых компонент.
  \item В процессе оптического полинга германосиликатного волокна излучением ближнего ИК-диапазона ширина кривых фазового синхронизма генерации второй гармоники монотонно уменьшается, а их пик сдвигается к расстройке, определяемой фазовой самомодуляцией и кросс-модуляцией. При этом выход эффективности ГВГ на постоянный уровень наступает до завершения эволюции кривых фазового синхронизма, поскольку указанные эффекты компенсируются одновременным ростом пиковой эффективности преобразования.
\end{enumerate}
%В папке Documents можно ознакомиться с решением совета из Томского~ГУ
%(в~файле \verb+Def_positions.pdf+), где обоснованно даются рекомендации
%по~формулировкам защищаемых положений.

{\reliability} полученных результатов обеспечивается высокой повторяемостью экспериментальных данных в пределах допустимой погрешности. Надёжность работы подтверждается применением современной приборной базы, а также сопровождением экспериментов численным моделированием. Основные результаты диссертации опубликованы в рецензируемых научных журналах.

\makeatletter
\setcounter{citeregistered}{0}
\makeatother

\ifnumequal{\value{bibliosel}}{0}
{%%% Встроенная реализация с загрузкой файла через движок bibtex8. (При желании, внутри можно использовать обычные ссылки, наподобие `\cite{vakbib1,vakbib2}`).
%    {\publications} 
    {\probation} Основные результаты по теме диссертации изложены
    в~XX~печатных изданиях,
    X из которых изданы в журналах, рекомендованных ВАК,
    X "--- в тезисах докладов.
    
    
    Общее число докладов на российских и международных конференциях~--- 16. Из них: 10 докладов на 8 международных конференциях (ICLO 2022, 2024; ALT 2025; LLPh 2023, 2025; Фотоника и информационная оптика 2022, 2023) и 6 докладов на 4 всероссийских конференциях (МФТИ 2022–2024; Самарская конкурс-конференция по оптике, лазерной физике и физике плазмы 2023).
}%
{%%% Реализация пакетом biblatex через движок biber
    \begin{refsection}[bl-author, bl-registered]
        % Это refsection=1.
        % Процитированные здесь работы:
        %  * подсчитываются, для автоматического составления фразы "Основные результаты ..."
        %  * попадают в авторскую библиографию, при usefootcite==0 и стиле `\insertbiblioauthor` или `\insertbiblioauthorgrouped`
        %  * нумеруются там в зависимости от порядка команд `\printbibliography` в этом разделе.
        %  * при использовании `\insertbiblioauthorgrouped`, порядок команд `\printbibliography` в нём должен быть тем же (см. biblio/biblatex.tex)
        %
        % Невидимый библиографический список для подсчёта количества публикаций:
        \phantom{\printbibliography[heading=nobibheading, section=1, env=countauthorvak,          keyword=biblioauthorvak]%
        \printbibliography[heading=nobibheading, section=1, env=countauthorwos,          keyword=biblioauthorwos]%
        \printbibliography[heading=nobibheading, section=1, env=countauthorscopus,       keyword=biblioauthorscopus]%
        \printbibliography[heading=nobibheading, section=1, env=countauthorconf,         keyword=biblioauthorconf]%
        \printbibliography[heading=nobibheading, section=1, env=countauthorother,        keyword=biblioauthorother]%
        \printbibliography[heading=nobibheading, section=1, env=countregistered,         keyword=biblioregistered]%
        \printbibliography[heading=nobibheading, section=1, env=countauthorpatent,       keyword=biblioauthorpatent]%
        \printbibliography[heading=nobibheading, section=1, env=countauthorprogram,      keyword=biblioauthorprogram]%
        \printbibliography[heading=nobibheading, section=1, env=countauthor,             keyword=biblioauthor]%
        \printbibliography[heading=nobibheading, section=1, env=countauthorvakscopuswos, filter=vakscopuswos]%
        \printbibliography[heading=nobibheading, section=1, env=countauthorscopuswos,    filter=scopuswos]}%
        %
        \nocite{*}%
        \nocite{ostapiv2025phase}
        %
        {\publications} Основные результаты по теме диссертации изложены в~\arabic{citeauthor}~печатных изданиях,
        \arabic{citeauthorvak} из которых изданы в журналах, рекомендованных ВАК, Web of~Science и Scopus%
%        \ifnum \value{citeauthorscopuswos}>0%
%            , \arabic{citeauthorscopuswos} "--- в~периодических научных журналах, индексируемых Web of~Science и Scopus%
%        \fi%
        \ifnum \value{citeauthorconf}>0%
            , \arabic{citeauthorconf} "--- в~тезисах докладов.
        \else%
            .
        \fi%
        \ifnum \value{citeregistered}=1%
            \ifnum \value{citeauthorpatent}=1%
                Зарегистрирован \arabic{citeauthorpatent} патент.
            \fi%
            \ifnum \value{citeauthorprogram}=1%
                Зарегистрирована \arabic{citeauthorprogram} программа для ЭВМ.
            \fi%
        \fi%
        \ifnum \value{citeregistered}>1%
            Зарегистрированы\ %
            \ifnum \value{citeauthorpatent}>0%
            \formbytotal{citeauthorpatent}{патент}{}{а}{}%
            \ifnum \value{citeauthorprogram}=0 . \else \ и~\fi%
            \fi%
            \ifnum \value{citeauthorprogram}>0%
            \formbytotal{citeauthorprogram}{программ}{а}{ы}{} для ЭВМ.
            \fi%
        \fi%
        % К публикациям, в которых излагаются основные научные результаты диссертации на соискание учёной
        % степени, в рецензируемых изданиях приравниваются патенты на изобретения, патенты (свидетельства) на
        % полезную модель, патенты на промышленный образец, патенты на селекционные достижения, свидетельства
        % на программу для электронных вычислительных машин, базу данных, топологию интегральных микросхем,
        % зарегистрированные в установленном порядке.(в ред. Постановления Правительства РФ от 21.04.2016 N 335)
    \end{refsection}%
    \begin{refsection}[bl-author, bl-registered]
        % Это refsection=2.
        % Процитированные здесь работы:
        %  * попадают в авторскую библиографию, при usefootcite==0 и стиле `\insertbiblioauthorimportant`.
        %  * ни на что не влияют в противном случае
        \nocite{vakbib2}%vak
        \nocite{patbib1}%patent
        \nocite{progbib1}%program
        \nocite{bib1}%other
        \nocite{confbib1}%conf
    \end{refsection}%
        %
        % Всё, что вне этих двух refsection, это refsection=0,
        %  * для диссертации - это нормальные ссылки, попадающие в обычную библиографию
        %  * для автореферата:
        %     * при usefootcite==0, ссылка корректно сработает только для источника из `external.bib`. Для своих работ --- напечатает "[0]" (и даже Warning не вылезет).
        %     * при usefootcite==1, ссылка сработает нормально. В авторской библиографии будут только процитированные в refsection=0 работы.
}


{\contribution} Все использованные в диссертации экспериментальные и теоретические результаты получены автором лично или при определяющем его участии. Материалы, представленные в работе, получены в результате экспериментальных исследований, выполненных автором на кафедре фотоники (базовая организация <<ООО ВПГ <<Лазеруан>>) физтех-школы фотоники электроники и молекулярной физики МФТИ (национальный исследовательский университет).

\textbf{Объем и структура работы.} Диссертация состоит из~введения,
\formbytotal{totalchapter}{глав}{ы}{}{} и
заключения.
%\formbytotal{totalappendix}{приложен}{ия}{ий}{}.
%% на случай ошибок оставляю исходный кусок на месте, закомментированным
%Полный объём диссертации составляет  \ref*{TotPages}~страницу
%с~\totalfigures{}~рисунками и~\totaltables{}~таблицами. Список литературы
%содержит \total{citenum}~наименований.
%
Полный объём диссертации составляет
\formbytotal{TotPages}{страниц}{у}{ы}{}, включая
\formbytotal{totalcount@figure}{рисун}{ок}{ка}{ков} и
\formbytotal{totalcount@table}{таблиц}{у}{ы}{}.
Список литературы содержит
\formbytotal{citenum}{наименован}{ие}{ия}{ий}.

%При использовании пакета \verb!biblatex! будут подсчитаны все работы, добавленные
%в файл \verb!biblio/author.bib!. Для правильного подсчёта работ в~различных
%системах цитирования требуется использовать поля:
%\begin{itemize}
%        \item \texttt{authorvak} если публикация индексирована ВАК,
%        \item \texttt{authorscopus} если публикация индексирована Scopus,
%        \item \texttt{authorwos} если публикация индексирована Web of Science,
%        \item \texttt{authorconf} для докладов конференций,
%        \item \texttt{authorpatent} для патентов,
%        \item \texttt{authorprogram} для зарегистрированных программ для ЭВМ,
%        \item \texttt{authorother} для других публикаций.
%\end{itemize}
%Для подсчёта используются счётчики:
%\begin{itemize}
%        \item \texttt{citeauthorvak} для работ, индексируемых ВАК,
%        \item \texttt{citeauthorscopus} для работ, индексируемых Scopus,
%        \item \texttt{citeauthorwos} для работ, индексируемых Web of Science,
%        \item \texttt{citeauthorvakscopuswos} для работ, индексируемых одной из трёх баз,
%        \item \texttt{citeauthorscopuswos} для работ, индексируемых Scopus или Web of~Science,
%        \item \texttt{citeauthorconf} для докладов на конференциях,
%        \item \texttt{citeauthorother} для остальных работ,
%        \item \texttt{citeauthorpatent} для патентов,
%        \item \texttt{citeauthorprogram} для зарегистрированных программ для ЭВМ,
%        \item \texttt{citeauthor} для суммарного количества работ.
%\end{itemize}
%% Счётчик \texttt{citeexternal} используется для подсчёта процитированных публикаций;
%% \texttt{citeregistered} "--- для подсчёта суммарного количества патентов и программ для ЭВМ.
%
%Для добавления в список публикаций автора работ, которые не были процитированы в
%автореферате, требуется их~перечислить с использованием команды \verb!\nocite! в
%\verb!Synopsis/content.tex!.
